% ============================================================
% Chapter 6: Experimental Design
%   Section: Model Benchmarks
%     Subsection: Verification Models
% Included from chapters/06-experimental-design/06-model-benchmarks/model-benchmarks.tex
% Source: docs/roadmap.md (Sprint 7; P1.4 of the thesis priorities);
% docs/decisions/evaluation.md, "Metrics" (cli table) and "Choosing the
% models" (2026-10-19).
% ============================================================

\subsection{Verification Models}
\label{subsec:benchmark-verification}

The verifier is compared across models on the custom set
(Table~\ref{tab:priorities}, P1): the local baseline,
\texttt{llama3.2:3b}; two openly licensed local models of the same
family at two sizes, \texttt{qwen3.5:4b} and \texttt{qwen3.5:9b}, with
their reasoning mode switched off; and, if time allows, one hosted model
reached through an OpenAI-compatible endpoint. Every other part of the
pipeline is held fixed --- thresholds, weights, engines, corpus
snapshot --- so the runs differ by the model alone, and each run's key
records that they do, the reasoning mode included: a hybrid model told
not to think and the same model thinking answer differently under one
name.

\paragraph{Chosen on one half, reported on the other.} Since the split
of Subsection~\ref{subsec:custom-split} exists, the models are compared on the
\emph{train} half, and the one chosen is then run once on the
\emph{test} half to report it; the test half chooses nothing. The rule
was written before any candidate's result was read: a candidate replaces
the baseline when its accuracy gain's paired 95\% interval over the
claims both runs scored lies above zero; when the interval includes
zero, the faster model is kept. With 69 claims a paired interval of
roughly ten points is expected, so only large differences separate
models. Two runs of the same model already differ --- the run-to-run
consistency of Subsection~\ref{subsec:consistency} and Table~\ref{tab:consistency} --- so a
small gain is not read as one.

\paragraph{What is set side by side.} Each model is one harness run.
The models are set side by side, one row each, with accuracy and
macro-$F_1$ and their intervals, coverage, selective accuracy, Cohen's
$\kappa$, invalid outputs, claims set apart, tokens per claim and per
second, the cost per hundred claims, and the \textbf{response time}: the
seconds a claim takes at the median and the 90th percentile, and the
model's own answer time per call. A run's report breaks the claim's time
into its stages --- the claim's entities, evidence retrieval, ranking,
the model's verdict, and a second look where one was taken --- with each
stage's share. The model's span includes waiting for its turn at the
model; the answer time does not; the gap between them is a busy model,
not a slow one. Response times compare only between runs made alone on
the machine, and each model is run alone.

The same models are repeated under the evidence conditions of RQ2 ---
no evidence, retrieved evidence, the annotator's evidence --- only for
the chosen one, so that the cost of the comparison stays within the
plan.
